\subsection{本性赋值的逼近}

由于 Krull 整环的本性赋值是离散且赋范的，任意两个本性赋值都不等价，因而它们是独立的（第 VI 章，§ 7，第 2 节）。因此可以对它们应用逼近定理的推论 2（同上，定理 1）：给定一些 \(n_i \in \mathbf{Z}\) 以及有限个互异的本性赋值 \(v_i\)，存在 \(x \in \mathbf{K}\)，使得 \(v_i(x) = n_i\) 对所有 \(i\) 都成立。不过，这里有一个更精确的结果：

命题 9. 设 \(v_1, \ldots, v_r\) 是 Krull 整环 A 的互异本性赋值，\(n_1, \ldots, n_r\) 是有理整数。则 A 的分式域 K 中存在元素 \(x\)，使得 \(v_i(x) = n_i\) 对 \(1 \leqslant i \leqslant r\) 都成立，并且 \(v(x) \geqslant 0\) 对 A 的每个本性赋值 \(v\)（不同于 \(v_1, \ldots, v_r\)）都成立。

设 \(\mathfrak{p}_1, \ldots, \mathfrak{p}_r\) 是分别对应于赋值 \(v_1, \ldots, v_r\) 的 A 的除性理想。存在 \(y \in K\)，使得 \(v_i(y) = n_i\) 对 \(1 \leqslant i \leqslant r\) 都成立（第 VI 章，§ 7，第 2 节，定理 1 的推论 1）。A 的本性赋值 \(w_1, \ldots, w_s\) 不同于 \(v_i\)，且 \(w_j(y) = -m_j\) 满足 \(<0\)，其数目是有限的；设相应的理想为 \(\mathfrak{q}_1, \ldots, \mathfrak{q}_s\)。\(\mathfrak{p}_1, \ldots, \mathfrak{p}_r, \mathfrak{q}_1, \ldots, \mathfrak{q}_s\) 之间不存在包含关系，因为这些理想对应于极端除子，而这些理想是素理想（命题 6 的推论 1）。因此，整理想 \(\mathfrak{a} = \mathfrak{q}_1^{m_1} \ldots \mathfrak{q}_s^{m_s}\) 不包含于任何一个 \(\mathfrak{p}_i\)（第 II 章，§ 1，第 1 节，命题 1），所以也不包含于它们的并中（同上，命题 2）。因此存在 \(z \in \mathfrak{a}\)，使得 \(z \notin \mathfrak{p}_i\) 对 \(1 \leqslant i \leqslant r\) 都成立；于是 \(v_1(z) = \cdots = v_r(z) = 0\)，且 \(w_j(z) \geqslant m_j\) 对 \(1 \leqslant j \leqslant s\) 都成立；从而元素 \(x = yz\) 解决了问题。

推论 1. 设 A 是 Krull 整环，K 是其分式域，a、b、c 是 A 的三个除性分式理想，且 a ⊂ b。则存在 x ∈ K，使得 a = b ∩ xc。

设 \((v_{\iota})_{\iota \in I}\) 是 A 的本性赋值族，设 \((m_{\iota})\)（分别地，\((n_{\iota})\)、\((p_{\iota})\)）是一个有理整数族（除有限个指标外其余值均为零），使得 \(\mathfrak{a}\)（分别地，\(\mathfrak{b}, \mathfrak{c}\)）是 \(x \in K\) 的集合，其中 \(v_{\iota}(x) \geqslant m_{\iota}\)（分别地，\(n_{\iota}, p_{\iota}\)）对所有 \(\iota \in I\) 都成立（命题 5，第 4 节）。集合 J 由指标 \(\iota \in I\) 中满足 \(m_{\iota} > n_{\iota}\) 的组成，且 J 是有限的。由于除有限个指标外 \(p_{\iota} = m_{\iota} = 0\)，命题 9 表明存在 \(x \in K^{*}\)，使得 \(v_{\iota}(x^{-1}) + m_{\iota} = p_{\iota}\) 对 \(\iota \in J\) 成立，并且

\[
v _ {\iota} (x ^ {- 1}) + m _ {\iota} \geqslant p _ {\iota}
\]

当 \(\iota \in \mathbf{I} - \mathbf{J}\) 时成立。于是对所有 \(\iota \in \mathbf{I}\)，\(m_{\iota} = \sup (n_{\iota}, v_{\iota}(x) + p_{\iota})\)。从而 \(\mathfrak{a} = \mathfrak{b} \cap x\mathfrak{c}\)。

推论 2. 设 A 是 Krull 整环。A 的分式理想 \(\mathfrak{a}\) 为除性的充要条件是它是两个分式主理想的交。

充分性显然（第 1 节，定义 2）。必要性由推论 1 得出：取 b 和 c 为主理想，并使 b ⊃ a。

